\documentclass[conference]{IEEEtran}
\IEEEoverridecommandlockouts

\usepackage{cite}
\usepackage{amsmath,amssymb,amsfonts}
\usepackage{graphicx}
\usepackage{textcomp}
\usepackage{xcolor}
\usepackage{url}

\begin{document}

\title{Multi Omics Integration}

\author{
\IEEEauthorblockN{Aman Kumar Singh}
\IEEEauthorblockA{
SEMTM0045 – Data Science Methods and Practice\\
University of Bristol}
}

\maketitle

\begin{abstract}
Provide a concise overview of the research problem, the identified gap, the proposed methodological direction, and the expected impact of the project.
\end{abstract}

\begin{IEEEkeywords}
keyword1, keyword2, ..
\end{IEEEkeywords}

\section{Introduction}
\footnote {GitHub link: \url{http:\\}}


\section{Literature Review}

Multi-omics integration has significantly advanced personalized medicine by allowing a more comprehensive understanding of complex biological systems beyond single-omics analyzes \cite{acharya_comprehensive_2024}. By integrating diverse modalities such as genomics, transcriptomics and epigenomics, these approaches offer deeper insight into disease mechanisms and patient heterogeneity \cite{hong_multi-omics_2025}. However, despite these advancements, the field continues to face fundamental computational, methodological, and evaluation challenges, most notably the curse of dimensionality \cite{guarino_intrinsic-dimension_2023}. Multi-omics datasets are inherently characterised by extremely high-dimensional feature spaces, coupled with relatively small sample sizes and low signal-to-noise ratios\cite{kioroglou_multi-omic_2025,hong_multi-omics_2025}. This high-dimensional, low sample size (HDLSS) setting necessitates rigorous preprocessing pipelines, including standardised normalisation, robust batch effect correction, data imputation, and effective dimensionality reduction to enable meaningful downstream analysis \cite{abdelaziz_multi-omics_2024}.\\
Within this context, two broad methodological clusters can be identified: (i) interpretable statistical and network-based approaches, and (ii) complex latent representation learning methods. Although dimensionality reduction is widely adopted to address high-dimensional multi-omics data, these methodological categories differ substantially in the extent to which they preserve biological interpretability. Several studies fall within the latter category and employ approaches that do not retain explicit feature-level representations\cite{shakyawar_icluf_nodate, hong_multi-omics_2025, yang_mdicc_2022}. For instance, Yang et al. \cite{yang_mdicc_2022} perform dimensionality reduction implicitly through MDICC via low-rank subspace learning, without constructing an explicit latent feature space, thereby limiting biological interpretability despite strong clustering performance. Similarly, Shakyawar et al. \cite{shakyawar_icluf_nodate} reduce dimensional complexity through iCluF via iterative refinement of patient–patient similarity networks, but provide limited insight into the underlying omics features contributing to the learned representations. Collectively, these methods reflect a broader trend in which improvements in clustering performance and computational efficiency are achieved at the cost of transparency and biological interpretability.\\
In contrast, classical approaches such as PCA  and NMF provide explicit feature contributions, enabling clearer biological interpretation \cite{downing_primer_nodate}. Network-based methods such as Similarity Network Fusion (SNF)\cite{cantini_benchmarking_2021} integrate multi-omics through fused similarity networks, while factor-based models such as MOFA \cite{darkaoui_integrative_2026} learn latent factors that retain partial biological interpretability. However, these methods are often limited in capturing complex nonlinear relationships inherent in biological systems. This contrast highlights a fundamental trade-off in the literature between interpretability and modelling expressiveness, with recent methods increasingly favouring predictive performance over biological transparency.\\
Beyond dimensionality reduction, batch effect correction represents an equally critical yet frequently underemphasised preprocessing step. Batch effects, arising from variations in experimental conditions, sequencing platforms, or laboratory protocols, can obscure true biological signals and introduce systematic bias into downstream analyses \cite{goh_why_2017}. Consequently, even well-designed dimensionality reduction techniques may fail when applied to uncorrected data. Despite this, several studies neglect robust batch correction strategies\cite{shakyawar_icluf_nodate, yang_mdicc_2022, hira_integrated_2021}. For example, Hira et al. \cite{hira_integrated_2021} integrate multi-source omics data (TCGA and GDC) without explicit batch correction; although normalisation and sample intersection are applied, their direct concatenation likely preserves technical variability, leading to representations influenced by batch-driven artefacts rather than biological signals.
Similarly, Yang et al. \cite{yang_mdicc_2022} focus on data harmonisation via MDICC but do not explicitly address cross-platform batch effects.
Zhang et al. \cite{zhang_integrative_2025} employ ComBat, while Kioroglou et al. \cite{kioroglou_multi-omic_2025} use covariate filtering, demonstrating improved control of non-biological variation. The omission of such steps constitutes a methodological gap, as uncorrected technical noise propagates through the analytical pipeline and compromises downstream validity.\\
At the downstream analysis stage, deep learning-based approaches constitute a second major methodological cluster in multi-omics research. Architectures such as ANNs, GNNs, GCNs, and transformer-based models are widely employed to capture complex nonlinear relationships and often outperform traditional machine learning methods in predictive tasks\cite{sartori_comprehensive_2025}. However, despite these performance gains, such models remain inherently black-box, offering limited feature-level interpretability and biological insight, thereby constraining their clinical applicability \cite{kang_roadmap_2022}. For example, Wu et al. \cite{wu_deepmoic_2024} demonstrate superior performance with DeepMoIC over classical classifiers such as SVM, Random Forest, and KNN, yet fail to provide clear attribution of the molecular features driving the model's predictions, limiting its biological relevance. Although post-hoc Explainable AI (XAI) frameworks, including SHAP, LIME, and Integrated Gradients, have been introduced to extract feature-level importance from such complex models, their adoption in multi-omics studies remains limited, and they provide only approximate explanations that do not fully resolve the inherent opacity of deep learning models \cite{noauthor_full_nodate}.\\
This limitation is further amplified by the HDLSS nature of multi-omics data. Deep learning models, which are inherently data-intensive, are prone to overfitting in such settings, leading to reduced generalisability. In contrast, traditional machine learning methods—particularly ensemble approaches such as XGBoost and Random Forest—demonstrate greater robustness and provide inherent feature importance measures that enhance interpretability. Empirical comparisons reinforce this distinction: in MOGONET [11], XGBoost outperforms the deep learning model on certain datasets and achieves comparable performance on others, while in [5], LightGBMXT achieves superior predictive performance (AUC = 0.9727) compared to multiple deep learning architectures. These findings challenge the implicit assumption that deep learning is universally superior, instead suggesting that simpler models may be more appropriate in HDLSS settings.

\section{Proposed Methodology}
\section{Data, Ethics, and Practical Considerations}
\section{Project plan}
\section{Non-Technical Summary}

\bibliographystyle{IEEEtran}
\bibliography{references}

\end{document}
